% Introduction section

Physics-Informed Neural Networks (PINNs) have emerged as a powerful paradigm for solving partial differential equations (PDEs) using deep learning~\cite{raissi2019physics}. The key idea is to embed physical constraints directly into the neural network training process through the loss function.

Recently, Taylor, Pardo, and Muga~\cite{taylor2023deep} introduced the Deep Fourier Residual (DFR) method, which uses the $H^{-1}$ dual norm of the weak residual as the loss function. This approach provides several advantages over traditional PINNs:
\begin{itemize}
    \item The loss is equivalent to the $H^1$ error for well-posed problems
    \item Reliable error estimation during training
    \item Applicability to problems with limited regularity
\end{itemize}

\subsection{Motivation}

% Describe the motivation for your follow-up work
% What limitations of DFR are you addressing?
% What new applications are you enabling?

[Describe the motivation for your extension here]

\subsection{Contributions}

The main contributions of this work are:
\begin{enumerate}
    \item [First contribution]
    \item [Second contribution]
    \item [Third contribution]
</enumerate>

\subsection{Organization}

The remainder of this paper is organized as follows. Section~\ref{sec:background} reviews the necessary background on PINNs and the DFR method. Section~\ref{sec:methodology} presents our proposed extensions. Section~\ref{sec:theory} provides theoretical analysis. Section~\ref{sec:experiments} describes our numerical experiments, and Section~\ref{sec:results} discusses the results. Finally, Section~\ref{sec:conclusion} concludes the paper.
